\subsection{6. File-by-File Reference: What to Write
Where}\label{file-by-file-reference-what-to-write-where}

\subsubsection{\texorpdfstring{\texttt{experiments/NNN\_folder\_name/README.md}}{experiments/NNN\_folder\_name/README.md}}\label{experimentsnnn_folder_namereadme.md}

This is the most important file in each experiment. Think of it as the
\textbf{scientific report} for that replication.

{\def\LTcaptype{none} % do not increment counter
\small
\begin{longtable}[]{@{}
  >{\raggedright\arraybackslash}p{(\linewidth - 2\tabcolsep) * \real{0.3913}}
  >{\raggedright\arraybackslash}p{(\linewidth - 2\tabcolsep) * \real{0.6087}}@{}}
\toprule\noalign{}
\begin{minipage}[b]{\linewidth}\raggedright
Section
\end{minipage} & \begin{minipage}[b]{\linewidth}\raggedright
What to write
\end{minipage} \\
\midrule\noalign{}
\endhead
\bottomrule\noalign{}
\endlastfoot
\textbf{Paper Metadata} & Fill in Authors, Year, Venue, DOI, Code URL,
language, ARA link, Excel row \\
\textbf{Replication Status} & Keep this badge accurate --- it drives the
master dashboard \\
\textbf{Objective} & List exactly which tables/figures you plan to
reproduce. Check them off when done. \\
\textbf{Environment Setup} & Write the exact commands to set up the
environment from scratch. Assume a reader starts with nothing. \\
\textbf{How to Run} & Write the exact commands to run the experiment. \\
\textbf{Key Results} & Fill in the comparison table. This is the
scientific output. \\
\textbf{Observations \& Deviations} & Be honest about differences from
the paper. \\
\end{longtable}
}

\subsubsection{\texorpdfstring{\texttt{experiments/NNN\_folder\_name/REPLICATION\_LOG.md}}{experiments/NNN\_folder\_name/REPLICATION\_LOG.md}}\label{experimentsnnn_folder_namereplication_log.md}

This is your \textbf{daily diary}. Write in it every time you work on
the experiment. New entries go at the \textbf{top}.

Format:

\begin{minted}{text}
## [YYYY-MM-DD] What You Did

**Time spent**: Xh Xm

### What was done
- ...

### What worked
- ...

### Issues encountered
- ...

### Next steps
- ...
\end{minted}

\subsubsection{\texorpdfstring{\texttt{experiments/NNN\_folder\_name/requirements.txt}}{experiments/NNN\_folder\_name/requirements.txt}}\label{experimentsnnn_folder_namerequirements.txt}

Contains the \textbf{Python packages} needed to run the experiment.
Update it after setting up the environment:

\begin{minted}{bash}
pip freeze > requirements.txt
\end{minted}

\subsubsection{\texorpdfstring{\texttt{experiments/NNN\_folder\_name/config/default.yaml}}{experiments/NNN\_folder\_name/config/default.yaml}}\label{experimentsnnn_folder_nameconfigdefault.yaml}

The \textbf{experiment configuration} file. Use it to track the key
parameters of your replication run: - Random seeds - Dataset paths -
Model hyperparameters - Evaluation settings

This gives you a single place to see what parameters produced your
results.

\subsubsection{\texorpdfstring{\texttt{shared/metrics/regression.py} and
\texttt{shared/metrics/classification.py}}{shared/metrics/regression.py and shared/metrics/classification.py}}\label{sharedmetricsregression.py-and-sharedmetricsclassification.py}

\textbf{Only use these if the original paper doesn't provide its own
metric implementations.} These are convenience implementations. The
paper's own metrics always take priority.
